\chapter{Error Analysis, Discussion, and Limitations}
\label{chap:discussion}
\updateheader{Chapter 7. Error Analysis, Discussion, and Limitations}

%==============================================================================
% 写作时间：Day 4。参照学长 Ch5：解释结果为什么是这样、机制实际起了什么作用、
% 哪些结论证据不足、局限性。
%==============================================================================

\section{Error Taxonomy}
\label{sec:error_taxonomy}
% 可沿用会议工作的错误类型：topic mismatch、over-generalization、weak grounding、
% answer-graph misalignment；新增：repair regression、unsupported relation 等。
\todo{Table of error types with counts and definitions.}

\section{Case Studies}
\label{sec:case_studies}
% 三个案例：成功；修复后成功（before/after）；失败。配 OPCloud 图（FIG03–FIG05）。
\todo{Write this section.}

\section{Interpretation of the Results}
\label{sec:interpretation}
\todo{Write this section.}

\section{Threats to Validity}
\label{sec:threats}
% 样本量小；单一模型；LLM judge 的偏差；评审者资质；
% 预训练数据可能包含公开的 MedQA 题目（数据污染）；关键词筛选可能漏掉部分心血管题。
\todo{Write this section.}

\section{Limitations}
\label{sec:limitations}
% 不评估临床正确性；OPM 表示是简化版；OPCloud 流程依赖人工。
\todo{Write this section.}
